% GENERATED FILE. Edit canonical lessons, then run npm run generate:books.
% canonical-source-hash: fnv1a64:bfba9ba25d638c4c
% canonical-lessons: PT-C516-orquestra, PT-C516-maestro, PT-C516-sinfonia, PT-C516-violoncelo, PT-C516-harpa

\chapter{A orquestra}
\label{ch:pt-a-orquestra}
\hlchaptermodality{516}

\begin{chapteropening}
\textbf{By the end of this chapter:} \emph{I can talk about an orchestra and its instruments.}

Complete the last lesson of chapter 516: I can talk about an orchestra and its instruments.
\end{chapteropening}

\section[a orquestra]{a orquestra — orchestra}

\label{lesson:PT-C516-orquestra}



\begin{quote}
Before the new one: say the Portuguese for to get injured, then the Portuguese for cheating.
\end{quote}

\subsection*{You'll want to know: a orquestra}
\textbf{a orquestra} — \textquotedblleft{}orchestra\textquotedblright{}.

\begin{quote}
A orquestra tocou durante duas horas.
\end{quote}

\begin{quote}
\emph{The orchestra played for two hours.}
\end{quote}

A musician in it is \emph{um músico da orquestra}.

\begin{cousinweb}
\emph{orquestra} comes from Latin \emph{orchēstra}, from Greek \emph{orkh\'{\={e}}stra} (EtymDB), the space in a theatre where the chorus danced, built on \emph{orkheîsthai} \textquotedblleft{}to dance\textquotedblright{}.
\end{cousinweb}

\begin{grammarlens}[title={What you've built}]
One more word for telling, explaining and describing at B1.
\end{grammarlens}

\subsection*{Your turn}
\begin{itemize}
\raggedright
  \item \emph{Say it:} \emph{a orquestra}
  \item \emph{Say it:} \emph{a orquestra}, once more
  \item \emph{Say it:} a sentence of your own with \emph{a orquestra}
  \item \emph{Recall:} say the Portuguese for to get injured, then the Portuguese for cheating, then say \emph{a orquestra} again
\end{itemize}

\begin{tcolorbox}[breakable,colback=teal!4,colframe=teal!35!black,title={Before you move on}]
What does a orquestra mean? (\textquotedblleft{}Orchestra\textquotedblright{}.) Say it once more.
\end{tcolorbox}
\section[o maestro]{o maestro — conductor}

\label{lesson:PT-C516-maestro}



\begin{quote}
Before the new one: say the Portuguese for cheating, then the Portuguese for orchestra.
\end{quote}

\subsection*{You'll want to know: o maestro}
\textbf{o maestro} — \textquotedblleft{}conductor\textquotedblright{}.

\begin{quote}
O maestro levantou os braços e a música começou.
\end{quote}

\begin{quote}
\emph{The conductor raised his arms and the music began.}
\end{quote}

Not a schoolteacher: in Portuguese, \emph{maestro} is the conductor.

\begin{cousinweb}
\emph{maestro} was borrowed from Italian \emph{maestro} \textquotedblleft{}master\textquotedblright{} (EtymDB), which is inherited from Latin \emph{magister} \textquotedblleft{}master, teacher\textquotedblright{}.
\end{cousinweb}

\begin{grammarlens}[title={What you've built}]
One more word for telling, explaining and describing at B1.
\end{grammarlens}

\subsection*{Your turn}
\begin{itemize}
\raggedright
  \item \emph{Say it:} \emph{o maestro}
  \item \emph{Say it:} \emph{o maestro}, once more
  \item \emph{Say it:} a sentence of your own with \emph{o maestro}
  \item \emph{Recall:} say the Portuguese for cheating, then the Portuguese for orchestra, then say \emph{o maestro} again
\end{itemize}

\begin{tcolorbox}[breakable,colback=teal!4,colframe=teal!35!black,title={Before you move on}]
What does o maestro mean? (\textquotedblleft{}Conductor\textquotedblright{}.) Say it once more.
\end{tcolorbox}
\section[a sinfonia]{a sinfonia — symphony}

\label{lesson:PT-C516-sinfonia}



\begin{quote}
Before the new one: say the Portuguese for orchestra, then the Portuguese for conductor.
\end{quote}

\subsection*{You'll want to know: a sinfonia}
\textbf{a sinfonia} — \textquotedblleft{}symphony\textquotedblright{}.

\begin{quote}
A sinfonia tem quatro partes.
\end{quote}

\begin{quote}
\emph{The symphony has four parts.}
\end{quote}

\emph{Orquestra sinfónica} — symphony orchestra.

\begin{cousinweb}
\emph{sinfonia} comes from Latin \emph{symphōnia} (EtymDB), from Greek \emph{symphōnía} \textquotedblleft{}a sounding together\textquotedblright{}, built on \emph{sýn} \textquotedblleft{}with\textquotedblright{} and \emph{phōn\'{\={e}}} \textquotedblleft{}sound\textquotedblright{}.
\end{cousinweb}

\begin{grammarlens}[title={What you've built}]
One more word for telling, explaining and describing at B1.
\end{grammarlens}

\subsection*{Your turn}
\begin{itemize}
\raggedright
  \item \emph{Say it:} \emph{a sinfonia}
  \item \emph{Say it:} \emph{a sinfonia}, once more
  \item \emph{Say it:} a sentence of your own with \emph{a sinfonia}
  \item \emph{Recall:} say the Portuguese for orchestra, then the Portuguese for conductor, then say \emph{a sinfonia} again
\end{itemize}

\begin{tcolorbox}[breakable,colback=teal!4,colframe=teal!35!black,title={Before you move on}]
What does a sinfonia mean? (\textquotedblleft{}Symphony\textquotedblright{}.) Say it once more.
\end{tcolorbox}
\section[o violoncelo]{o violoncelo — cello}

\label{lesson:PT-C516-violoncelo}



\begin{quote}
Before the new one: say the Portuguese for conductor, then the Portuguese for symphony.
\end{quote}

\subsection*{You'll want to know: o violoncelo}
\textbf{o violoncelo} — \textquotedblleft{}cello\textquotedblright{}.

\begin{quote}
O violoncelo tem um som grave e bonito.
\end{quote}

\begin{quote}
\emph{The cello has a deep, beautiful sound.}
\end{quote}

The player is \emph{o violoncelista}.

\begin{cousinweb}
\emph{violoncelo} comes from Italian \emph{violoncello} (EtymDB), a diminutive of \emph{violone}, a large viol.
\end{cousinweb}

\begin{grammarlens}[title={What you've built}]
One more word for telling, explaining and describing at B1.
\end{grammarlens}

\subsection*{Your turn}
\begin{itemize}
\raggedright
  \item \emph{Say it:} \emph{o violoncelo}
  \item \emph{Say it:} \emph{o violoncelo}, once more
  \item \emph{Say it:} a sentence of your own with \emph{o violoncelo}
  \item \emph{Recall:} say the Portuguese for conductor, then the Portuguese for symphony, then say \emph{o violoncelo} again
\end{itemize}

\begin{tcolorbox}[breakable,colback=teal!4,colframe=teal!35!black,title={Before you move on}]
What does o violoncelo mean? (\textquotedblleft{}Cello\textquotedblright{}.) Say it once more.
\end{tcolorbox}
\section[a harpa]{a harpa — harp}

\label{lesson:PT-C516-harpa}



\begin{quote}
Before the new one: say the Portuguese for symphony, then the Portuguese for cello.
\end{quote}

\subsection*{You'll want to know: a harpa}
\textbf{a harpa} — \textquotedblleft{}harp\textquotedblright{}.

\begin{quote}
A harpa é um dos instrumentos mais antigos.
\end{quote}

\begin{quote}
\emph{The harp is one of the oldest instruments.}
\end{quote}

The player is \emph{o harpista} or \emph{a harpista}.

\begin{cousinweb}
\emph{harpa} goes back, through Old Portuguese \emph{arpa}, to Late Latin \emph{harpa}, which was taken from Frankish \emph{harpa} (EtymDB).
\end{cousinweb}

\begin{grammarlens}[title={What you've built}]
One more word for telling, explaining and describing at B1.
\end{grammarlens}

\subsection*{Your turn}
\begin{itemize}
\raggedright
  \item \emph{Say it:} \emph{a harpa}
  \item \emph{Say it:} \emph{a harpa}, once more
  \item \emph{Say it:} a sentence of your own with \emph{a harpa}
  \item \emph{Recall:} say the Portuguese for symphony, then the Portuguese for cello, then say \emph{a harpa} again
\end{itemize}

\begin{tcolorbox}[breakable,colback=teal!4,colframe=teal!35!black,title={Before you move on}]
What does a harpa mean? (\textquotedblleft{}Harp\textquotedblright{}.) Say it once more.
\end{tcolorbox}
